%! Matrices and Column Spaces — Unit 1, Lesson 1.3 — Group Activity (Answer Key)
\documentclass[10pt]{article}
\usepackage{linalg-article}
\usepackage{linalg-key}   % pulls in -boxes; do NOT also load -boxes
\newcommand{\TallMath}[1]{$\displaystyle #1\rule[-1.4em]{0pt}{3.2em}$}
\newcommand{\xx}{\mathbf{x}}
\newcommand{\bb}{\mathbf{b}}

\begin{document}

\pageheader{Unit 1, Lesson 1.3}{Group Activity --- Answer Key}

\vspace{0.10in}

\begin{headlinebox}{blush}
\small\textbf{Building Orders from Two Boxes.} A gift shop sells two boxes. Each box is a vector
$\begin{bsmallmatrix} \text{mugs} \\ \text{shirts} \end{bsmallmatrix}$, and the shop's matrix lines
them up as columns:
$A=\begin{bsmallmatrix} 2 & 1 \\ 1 & 3 \end{bsmallmatrix}$ ---
Classic $=\begin{bsmallmatrix} 2\\1 \end{bsmallmatrix}$, Deluxe $=\begin{bsmallmatrix} 1\\3 \end{bsmallmatrix}$.
If $\xx=\begin{bsmallmatrix} \text{\# Classic}\\ \text{\# Deluxe} \end{bsmallmatrix}$, then $A\xx$ is the
\emph{total order}. Work with your partner and \emph{show your steps}.
\end{headlinebox}

\vspace{0.10in}

\begin{tcolorbox}[colback=white, colframe=black!40, title=\textbf{Tier R --- Remediate},
    fonttitle=\bfseries, arc=1.5mm, left=3mm, right=3mm, top=2mm, bottom=2mm, breakable]
Columns of $A=\begin{bmatrix} 2 & 1 \\ 1 & 3 \end{bmatrix}$: Classic $=\begin{bsmallmatrix} 2\\1 \end{bsmallmatrix}$,
Deluxe $=\begin{bsmallmatrix} 1\\3 \end{bsmallmatrix}$.
\begin{enumerate}[leftmargin=1.7em, itemsep=8pt]
  \item Write out the two columns of $A$.
\begin{work}
  \mathbf{a}_1 &= \begin{bsmallmatrix} 2\\1 \end{bsmallmatrix} \\
  \mathbf{a}_2 &= \begin{bsmallmatrix} 1\\3 \end{bsmallmatrix}
\end{work}
  \item Compute $A\begin{bmatrix} 2 \\ 1 \end{bmatrix}$ as a combination of the columns
        (2 Classic, 1 Deluxe).
\begin{work}
  A\begin{bsmallmatrix} 2\\1 \end{bsmallmatrix}
      &= 2\begin{bsmallmatrix} 2\\1 \end{bsmallmatrix} + 1\begin{bsmallmatrix} 1\\3 \end{bsmallmatrix} \\
      &= \begin{bsmallmatrix} 4\\2 \end{bsmallmatrix} + \begin{bsmallmatrix} 1\\3 \end{bsmallmatrix} \\
      &= \begin{bsmallmatrix} 5\\5 \end{bsmallmatrix}
\end{work}
  \item Compute $A\begin{bmatrix} 3 \\ 0 \end{bmatrix}$. What does the answer mean in mugs and shirts?
\begin{work}
  A\begin{bsmallmatrix} 3\\0 \end{bsmallmatrix}
      &= 3\begin{bsmallmatrix} 2\\1 \end{bsmallmatrix} + 0\begin{bsmallmatrix} 1\\3 \end{bsmallmatrix} \\
      &= \begin{bsmallmatrix} 6\\3 \end{bsmallmatrix}
\end{work}
        \ansline{3 Classic boxes give 6 mugs and 3 shirts.}
\end{enumerate}
\end{tcolorbox}

\begin{tcolorbox}[colback=white, colframe=black!40, title=\textbf{Tier A --- Approaching Proficiency},
    fonttitle=\bfseries, arc=1.5mm, left=3mm, right=3mm, top=2mm, bottom=2mm, breakable]
Same shop, $A=\begin{bmatrix} 2 & 1 \\ 1 & 3 \end{bmatrix}$.
\begin{enumerate}[leftmargin=1.7em, itemsep=8pt]
  \item A customer wants a total order of $\begin{bmatrix} 5 \\ 5 \end{bmatrix}$ (5 mugs, 5 shirts).
        Find $\xx$ with $A\xx=\begin{bsmallmatrix} 5\\5 \end{bsmallmatrix}$. (Hint: check Tier~R.)
\begin{work}
  2\begin{bsmallmatrix} 2\\1 \end{bsmallmatrix} + 1\begin{bsmallmatrix} 1\\3 \end{bsmallmatrix}
      &= \begin{bsmallmatrix} 5\\5 \end{bsmallmatrix} \\
  \xx &= \begin{bsmallmatrix} 2\\1 \end{bsmallmatrix}
      \ \text{(2 Classic, 1 Deluxe)}
\end{work}
  \item Are the two columns of $A$ parallel (multiples of each other)?
        \ansline{No --- $\begin{bsmallmatrix} 1\\3 \end{bsmallmatrix}$ is not a multiple of $\begin{bsmallmatrix} 2\\1 \end{bsmallmatrix}$. They are independent.}
  \item \textbf{Interpret.} Because the two boxes are \emph{independent}, is the column space a
        line or the whole plane? What does that say about which orders the shop can fill?
        \ansline{The whole plane --- two independent columns span every point. With the right (possibly fractional) number of each box, the shop can fill \emph{any} mug/shirt target.}
\end{enumerate}
\end{tcolorbox}

\begin{tcolorbox}[colback=white, colframe=black!40, title=\textbf{Tier E --- Extension},
    fonttitle=\bfseries, arc=1.5mm, left=3mm, right=3mm, top=2mm, bottom=2mm, breakable]
A different shop stocks two boxes that are \emph{parallel}:
$M=\begin{bmatrix} 2 & 4 \\ 1 & 2 \end{bmatrix}$, columns $\begin{bsmallmatrix} 2\\1 \end{bsmallmatrix}$ and
$\begin{bsmallmatrix} 4\\2 \end{bsmallmatrix}$.
\begin{enumerate}[leftmargin=1.7em, itemsep=8pt]
  \item Show that column~2 is a multiple of column~1. What line is the column space $C(M)$?
\begin{work}
  \begin{bsmallmatrix} 4\\2 \end{bsmallmatrix}
      &= 2\begin{bsmallmatrix} 2\\1 \end{bsmallmatrix}
\end{work}
        \ansline{$C(M)$ is the line of all multiples of $\begin{bsmallmatrix} 2\\1 \end{bsmallmatrix}$ --- points with twice as many mugs as shirts.}
  \item Can this shop fill an order of $\begin{bmatrix} 3 \\ 3 \end{bmatrix}$? Decide, and explain
        using the columns.
        \ansline{No. $\begin{bsmallmatrix} 3\\3 \end{bsmallmatrix}$ is not a multiple of $\begin{bsmallmatrix} 2\\1 \end{bsmallmatrix}$ (would need $3=2t$ and $3=t$ at once), so it is not on the line $C(M)$.}
  \item \textbf{Justify.} Explain why \emph{any} order this shop fills must have
        twice as many mugs as shirts --- no matter how many of each box it uses.
\begin{work}
  a\begin{bsmallmatrix} 2\\1 \end{bsmallmatrix} + b\begin{bsmallmatrix} 4\\2 \end{bsmallmatrix}
      &= \begin{bsmallmatrix} 2a+4b\\a+2b \end{bsmallmatrix} \\
  2a+4b &= 2(a+2b)
\end{work}
        \ansline{The top entry is always $2\times$ the bottom, so mugs $=2\times$ shirts for every order.}
\end{enumerate}
\end{tcolorbox}

\end{document}
